% TikZ 数据预处理流程图（Nature 风格）
% 配色: Wong (Nature Methods, 2011) 色盲安全调色板
% 布局: 水平流水线 + 纵向分支，简洁克制
\begin{tikzpicture}[
    font=\small,
    node distance=3mm and 6mm,
]

% Wong 调色板
\definecolor{wongBlue}{HTML}{0072B2}
\definecolor{wongSkyBlue}{HTML}{56B4E9}
\definecolor{wongGreen}{HTML}{009E73}
\definecolor{wongOrange}{HTML}{E69F00}
\definecolor{wongVermillion}{HTML}{D55E00}
\definecolor{wongRedPurple}{HTML}{CC79A7}

% 淡色填充
\definecolor{tintBlue}{HTML}{DCE8F4}
\definecolor{tintGreen}{HTML}{DFF0E8}
\definecolor{tintOrange}{HTML}{FDF2DC}
\definecolor{tintVermillion}{HTML}{F8E5D6}
\definecolor{tintSky}{HTML}{E4F0F9}
\definecolor{tintRP}{HTML}{F3E4ED}

% 节点样式
\tikzset{
    pipenode/.style={
        rounded corners=2pt, draw=#1, line width=0.6pt, fill=#1!8,
        minimum width=28mm, minimum height=10mm,
        align=center, inner sep=2.5pt, font=\small
    },
    pipenode/.default=wongBlue,
    iobox/.style={
        rounded corners=2pt, draw=#1, line width=0.5pt, fill=#1!6,
        minimum width=24mm, minimum height=8mm,
        align=center, inner sep=2pt, font=\footnotesize
    },
    iobox/.default=wongSkyBlue,
    arrow/.style={
        -{Latex[length=1.8mm, width=1.6mm]}, line width=0.55pt, black!70
    },
    phaselabel/.style={
        font=\footnotesize\bfseries, text=black!70
    },
}

% ---- 输入源（顶部两个） ----
\node[iobox=wongSkyBlue] (csv) {CSV 索引表};
\node[iobox=wongSkyBlue, right=28mm of csv] (struct) {PDB / SDF\\结构文件};

% ---- Stage 1: 索引读取与过滤 ----
\node[pipenode=wongBlue, below=8mm of $(csv)!0.5!(struct)$] (s1)
    {索引读取与过滤\\[-1pt]{\footnotesize \texttt{filter\_df}}};

% ---- Stage 2: 结构批量加载 ----
\node[pipenode=wongBlue, below=8mm of s1] (s2)
    {结构批量加载\\[-1pt]{\footnotesize \texttt{read\_pocket} · \texttt{sdf\_load}}};

% ---- Stage 3: 复合体构建 ----
\node[pipenode=wongGreen, below=8mm of s2] (s3)
    {复合体构建\\[-1pt]{\footnotesize \texttt{merge\_sdf\_pdb}}};

% ---- Stage 4: 图特征转换 ----
\node[pipenode=wongGreen, below=8mm of s3] (s4)
    {图特征转换\\[-1pt]{\footnotesize \texttt{complex\_to\_data} · \texttt{pmap\_chunked}}};

% ---- Stage 5: 有效性校验与标签回填 ----
\node[pipenode=wongOrange, below=8mm of s4] (s5)
    {有效性校验\\[-1pt]{\footnotesize \texttt{assign\_label2x}}};

% ---- Stage 6: 数据持久化 ----
\node[pipenode=wongOrange, below=8mm of s5] (s6)
    {LMDB / CSV 缓存写入\\[-1pt]{\footnotesize \texttt{LmdbPPI.make\_shared}}};

% ---- Stage 7: 时间维划分 ----
\node[pipenode=wongVermillion, below=8mm of s6] (s7)
    {时间维划分\\[-1pt]{\footnotesize \texttt{diffdock\_splits}}};

% ---- 输出（底部） ----
\node[iobox=wongRedPurple, below left=8mm and 2mm of s7] (train) {train 子集};
\node[iobox=wongRedPurple, below=8mm of s7] (valid) {valid 子集};
\node[iobox=wongRedPurple, below right=8mm and 2mm of s7] (test) {test 子集};

% ---- 箭头：输入→Stage1 ----
\draw[arrow] (csv.south)   -- (s1.north west);
\draw[arrow] (struct.south) -- (s1.north east);

% ---- 箭头：Stage间 ----
\draw[arrow] (s1) -- (s2);
\draw[arrow] (s2) -- (s3);
\draw[arrow] (s3) -- (s4);
\draw[arrow] (s4) -- (s5);
\draw[arrow] (s5) -- (s6);
\draw[arrow] (s6) -- (s7);

% ---- 箭头：Stage7→输出 ----
\draw[arrow] (s7.south) -- ++(0,-2mm) -| (train.north);
\draw[arrow] (s7.south) -- (valid.north);
\draw[arrow] (s7.south) -- ++(0,-2mm) -| (test.north);

% ---- 右侧阶段标签 ----
\node[phaselabel] at ([xshift=14mm]$(s1.south east)!0.5!(s2.north east)$) {数据准备};
\node[phaselabel] at ([xshift=14mm]$(s3.south east)!0.5!(s4.north east)$) {特征构建};
\node[phaselabel] at ([xshift=14mm]$(s5.south east)!0.5!(s6.north east)$) {质量控制};
\node[phaselabel] at ([xshift=14mm]s7.east) {评测划分};

\end{tikzpicture}
